\documentclass[11pt]{amsart}

\usepackage[T1]{fontenc}
\usepackage{amsmath,amssymb,amsthm,mathtools}
\usepackage[expansion=false]{microtype}
\usepackage[hidelinks]{hyperref}
\hypersetup{
  pdftitle={Divisibility of Concatenated Arithmetic Progressions and Champernowne Prefixes},
  pdfauthor={Qichao Wang and Daoyu Dong}
}

\newtheorem{theorem}{Theorem}[section]
\newtheorem{lemma}[theorem]{Lemma}
\newtheorem{corollary}[theorem]{Corollary}
\newtheorem{proposition}[theorem]{Proposition}
\theoremstyle{remark}
\newtheorem{remark}[theorem]{Remark}

\newcommand{\concat}{\mathbin{\Vert}}
\newcommand{\ord}{\operatorname{ord}}

\title[Divisibility of concatenations]{Divisibility of Concatenated Arithmetic Progressions and Champernowne Prefixes}
\author{Qichao Wang\textsuperscript{*}}
\author{Daoyu Dong\textsuperscript{*}}
\thanks{\textsuperscript{*}These authors contributed equally to this work.}
\date{September 10, 2026}

\begin{document}

\begin{abstract}
We study divisibility questions for integers obtained by concatenation in positional notation.
For fixed digit length, concatenating an arithmetic progression is an affine function of its
initial term, which yields an exact congruence criterion.  We use this to prove a global
existence theorem in every even base; in decimal this resolves the existence conjectures
recorded in OEIS A075001 and A075000.

We then consider the Champernowne prefixes
\[
P(k)=1\concat2\concat\cdots\concat k.
\]
Our main theorem proves that for every positive integer \(n\) there exists \(k\) such that
\(n\mid P(k)\), thereby resolving the well-definedness conjecture in OEIS A075002.  The
proof has two local ingredients.  Away from \(3\), fixed-digit blocks admit unit translations
and explicit state periods.  At \(3\), where the unit argument degenerates because
\(10^d\equiv1\pmod 3\), a separate block-shift congruence gives a \(3\)-adic lifting
mechanism.  The local roots are synchronized at one common prefix index by processing odd
prime factors in increasing order; every completed state period has no prime divisor larger
than the prime just processed.
\end{abstract}

\maketitle

\section{Introduction}

Let
\[
P(k)=1\concat2\concat\cdots\concat k
\]
denote the decimal integer obtained by concatenating \(1,2,\ldots,k\), and for \(n\ge1\) let
\[
C_n(k)=k\concat(k+1)\concat\cdots\concat(k+n-1).
\]
OEIS A075001 asks for the smallest \(k\) such that \(n\mid C_n(k)\), with value \(0\)
if no such \(k\) exists, while OEIS A075000 records the corresponding smallest quotient
\cite{oeis75001,oeis75000}.  OEIS A075002 asks for the smallest \(k\) such that
\(n\mid P(k)\), where \(P\) is also the sequence in OEIS A007908
\cite{oeis75002,oeis77908}.  On September 10, 2026, the three entries still recorded
these existence statements as conjectures.  We are not aware of a previous proof of the
full A075002 assertion; this paper makes no claim of priority.

Our first result is a general even-base construction for concatenated arithmetic progressions.
The second and main result is the complete existence theorem for Champernowne prefixes.
The proof strategy is modular.  We first choose a single sufficiently long digit block.  The
endpoint is made divisible by the factor of the modulus whose prime divisors divide the base, and the exceptional
\(3\)-power is installed by a three-way block shift.  The remaining odd prime powers are then
installed in increasing order, using state periods to preserve all earlier congruences.  The
fact that a completed state period has no prime divisor larger than the prime already processed
ensures that the next translation remains a unit; a finite capacity bound keeps every chosen
endpoint in the same digit block.

Concatenation of arithmetic progressions in arbitrary bases has also been studied by Luca and
Teguia Tabuguia, who derive explicit formulae and prove non-holonomicity
\cite{luca-teguia}.  Our emphasis is different: divisibility by an externally prescribed
modulus and, in particular, the three OEIS existence problems above.

\section{Fixed-digit arithmetic progressions}
\label{sec:fixed-digit-ap}

Let \(\beta\ge2\), let \(n\ge1\), let \(s\in\mathbb Z\), and let \(d\ge1\).  Put
\(B=\beta^d\).  Suppose that
\[
k,\ k+s,\ \ldots,\ k+(n-1)s
\]
are positive \(d\)-digit base-\(\beta\) integers.  Their concatenation is
\begin{equation}\label{eq:ap-concat}
C_{\beta,n,s,d}(k)
  =\sum_{j=0}^{n-1}(k+js)B^{\,n-1-j}.
\end{equation}
Define
\[
R_n(B)=\sum_{i=0}^{n-1}B^i,
\qquad
T_n(B)=\sum_{j=0}^{n-1}jB^{\,n-1-j}.
\]
Then
\begin{equation}\label{eq:ap-affine}
C_{\beta,n,s,d}(k)=kR_n(B)+sT_n(B).
\end{equation}

\begin{proposition}[Exact fixed-digit criterion]\label{prop:fixed-digit-criterion}
Let \(g=\gcd(R_n(B),n)\).  The congruence
\[
n\mid C_{\beta,n,s,d}(k)
\]
is solvable in \(k\) if and only if
\[
g\mid sT_n(B).
\]
When solvable, the solutions form one residue class modulo \(n/g\).
\end{proposition}

\begin{proof}
By \eqref{eq:ap-affine}, the desired divisibility is the linear congruence
\[
kR_n(B)\equiv-sT_n(B)\pmod n.
\]
The standard solvability criterion for a linear congruence gives the claim.  After division
by \(g\), the coefficient \(R_n(B)/g\) is invertible modulo \(n/g\).
\end{proof}

\begin{theorem}[Even-base AP divisibility]\label{thm:even-base-ap}
Let \(\beta\ge2\) be even.  For every \(n\ge1\) and every \(s\in\mathbb Z\), there
exist \(d\ge1\) and \(k\) such that all
\[
k,\ k+s,\ldots,k+(n-1)s
\]
are positive \(d\)-digit base-\(\beta\) integers and
\[
n\mid C_{\beta,n,s,d}(k).
\]
\end{theorem}

\begin{proof}
Write
\begin{equation}\label{eq:ap-factorization}
n=qm,
\end{equation}
where \(q\) is the largest divisor of \(n\) all of whose prime divisors divide \(\beta\).  Thus
\((q,m)=1\) and \((m,\beta)=1\).

Choose \(d\) arbitrarily large subject to
\[
q\mid\beta^d,
\qquad
\beta^d\equiv1\pmod m.
\]
Such \(d\) exist: the first condition holds for all sufficiently large \(d\), and the
second holds whenever \(d\) is a multiple of \(\ord_m(\beta)\), with no order condition
when \(m=1\).

Put \(B=\beta^d\), and let
\[
u_- =\min\{0,(n-1)s\},
\qquad
u_+ =\max\{0,(n-1)s\}.
\]
All blocks have \(d\) digits whenever
\begin{equation}\label{eq:ap-window}
\beta^{d-1}-u_-
\le k\le
\beta^d-1-u_+.
\end{equation}
The length of this interval tends to infinity with \(d\), so we may choose \(k\) in it such
that
\[
k+(n-1)s\equiv0\pmod q.
\]
Modulo \(q\), every term of the sum in \eqref{eq:ap-concat} except the final block is
multiplied by a positive power of \(B\).  Hence
\begin{equation}\label{eq:ap-smooth}
C_{\beta,n,s,d}(k)\equiv k+(n-1)s\equiv0\pmod q.
\end{equation}
Modulo \(m\), since \(B\equiv1\),
\begin{equation}\label{eq:ap-coprime}
C_{\beta,n,s,d}(k)
 \equiv nk+s\frac{n(n-1)}2\pmod m.
\end{equation}
The first term is divisible by \(m\).  If \(n\) is odd, then
\[
s\frac{n(n-1)}2=mqs\frac{n-1}{2}.
\]
If \(n\) is even, then \(2\mid\beta\), so the full \(2\)-part of \(n\) lies in \(q\);
hence \(q/2\in\mathbb Z\), and
\[
s\frac{n(n-1)}2=m\frac q2\,s(n-1).
\]
Thus the concatenation is \(0\) modulo both coprime factors \(q\) and \(m\), hence modulo
\(n\).
\end{proof}

\begin{corollary}[A075001]\label{cor:A075001}
For every \(n\ge1\) there exists a positive integer \(k\) such that
\[
n\mid k\concat(k+1)\concat\cdots\concat(k+n-1).
\]
Hence OEIS A075001 has no zero terms.
\end{corollary}

\begin{proof}
Take \(\beta=10\) and \(s=1\) in Theorem~\ref{thm:even-base-ap}.
\end{proof}

\begin{corollary}[A075000]\label{cor:A075000}
For every \(n\ge1\), OEIS A075000 is nonzero.  More precisely,
\[
A075000(n)=\frac{C_n(A075001(n))}{n}.
\]
\end{corollary}

\begin{proof}
Existence follows from Corollary~\ref{cor:A075001}.  It remains only to justify minimality.
For fixed \(n\), if \(k_2>k_1\), then comparison of the decimal strings defining
\(C_n(k_2)\) and \(C_n(k_1)\) shows that
\[
C_n(k_2)>C_n(k_1),
\]
including when one of the indices crosses a power of ten.  Thus the least divisible
concatenation occurs at the least witnessing index \(A075001(n)\).  Division by the fixed
positive integer \(n\) preserves the ordering.
\end{proof}

\section{Fixed-block identities for Champernowne prefixes}
\label{sec:fixed-block-identities}

Fix a digit length \(d\) and put \(B=10^d\).  Suppose all integers appended after an
endpoint \(k\) remain in this \(d\)-digit block.  For \(L\ge0\),
\begin{equation}\label{eq:block-transition}
P(k+L)
 =B^LP(k)+\sum_{j=1}^L(k+j)B^{L-j}.
\end{equation}
Set
\[
S_L=\sum_{r=0}^{L-1}B^r,
\qquad
T_L=\sum_{j=1}^LjB^{L-j}.
\]
Then
\begin{equation}\label{eq:shift-identity}
P(k+L)-P(k)
 =(B^L-1)P(k)+kS_L+T_L.
\end{equation}
A direct telescoping calculation gives
\begin{equation}\label{eq:geometric-T}
(B-1)T_L=BS_L-L.
\end{equation}

\section{Local lifting and state periods}
\label{sec:local-lifting}

\begin{lemma}[Unit translation]\label{lem:unit-translation}
Let \(p\) be an odd prime with \(p\nmid B(B-1)\), let
\[
h=\ord_p(B),
\qquad
U=T_h,
\]
and let \(a\ge0\).  If \(p\nmid c\) and \(L=chp^a\), then
\begin{equation}\label{eq:unit-translation}
P(k+L)-P(k)\equiv cp^aU\pmod{p^{a+1}},
\end{equation}
provided all appended integers remain in the same digit block.  Moreover \(p\nmid U\).
\end{lemma}

\begin{proof}
Modulo \(p\), \(B^h=1\) and \(S_h=0\).  From \eqref{eq:geometric-T},
\[
(B-1)U\equiv-h\pmod p.
\]
Since \(h\mid p-1\), both \(h\) and \(B-1\) are units modulo \(p\), hence so is \(U\).

By the lifting-the-exponent lemma,
\[
B^{chp^a}\equiv1\pmod{p^{a+1}}.
\]
Since \(B-1\) is a unit, this also gives
\[
S_L\equiv0\pmod{p^{a+1}}.
\]
Equation \eqref{eq:geometric-T} therefore yields
\[
T_L\equiv-\frac{chp^a}{B-1}\equiv cp^aU
\pmod{p^{a+1}}.
\]
Substitution into \eqref{eq:shift-identity} proves \eqref{eq:unit-translation}.
\end{proof}

\begin{lemma}[State period]\label{lem:state-period}
Under the hypotheses of Lemma~\ref{lem:unit-translation}, for every \(e\ge1\),
\begin{equation}\label{eq:state-period-length}
H=p^e\ord_{p^e}(B)
\end{equation}
is a period of the endpoint state modulo \(p^e\):
\begin{equation}\label{eq:state-period}
P(k+H)\equiv P(k)\pmod{p^e},
\qquad
k+H\equiv k\pmod{p^e}.
\end{equation}
Every multiple of \(H\) is also such a period.  Every prime divisor of \(H\) is at most
\(p\).
\end{lemma}

\begin{proof}
Let \(r=\ord_{p^e}(B)\).  Since \(r\mid H\), \(B^H\equiv1\pmod{p^e}\).  Because
\(B-1\) is a unit, \(S_H\equiv0\pmod{p^e}\).  Also \(p^e\mid H\), so
\eqref{eq:geometric-T} gives \(T_H\equiv0\pmod{p^e}\).  Equation
\eqref{eq:shift-identity} proves the first congruence in \eqref{eq:state-period}, and
\(p^e\mid H\) proves the second.  Iterating the state transition shows that every multiple
of \(H\) is again a period.

Finally,
\[
\ord_{p^e}(B)\mid\varphi(p^e)=p^{e-1}(p-1),
\]
so every prime divisor of \(H\) is at most \(p\).
\end{proof}

\begin{lemma}[Exceptional \(3\)-adic block shift]\label{lem:3-adic-shift}
Let \(B=10^d\), and suppose \(v_3(L)=a\ge1\).  Then
\begin{equation}\label{eq:3-adic-block-shift}
P(k+L)-P(k)
\equiv Lk+\frac{L(L+1)}2
\pmod{3^{a+1}}
\end{equation}
whenever the appended integers remain in one digit block.
\end{lemma}

\begin{proof}
Write \(B=1+c\).  Since \(B=10^d\), we have \(9\mid c\).  Expansion of the geometric
sum gives
\[
S_L-L=\sum_{s\ge1}c^s\binom{L}{s+1}.
\]
Using
\[
v_3\binom{L}{m}\ge a-v_3(m)
\]
and \(v_3(c)\ge2\), every term on the right has valuation at least \(a+1\).  Hence
\[
S_L\equiv L\pmod{3^{a+1}}.
\]

Similarly,
\[
T_L-\frac{L(L+1)}2
 =\sum_{s\ge1}c^s\binom{L+1}{s+2}.
\]
From
\[
\binom{L+1}{m}
 =\frac{L(L+1)}{m(m-1)}\binom{L-1}{m-2}
\]
we obtain
\[
v_3\!\left(c^s\binom{L+1}{s+2}\right)
 \ge 2s+a-v_3((s+1)(s+2)).
\]
For \(s\ge1\),
\[
v_3((s+1)(s+2))\le 2s-1;
\]
for \(s=1\) this is immediate, and for \(s\ge2\) it follows from
\(3^{2s}>(s+1)(s+2)\).  Thus every correction term is divisible by \(3^{a+1}\).

Finally LTE gives \(B^L\equiv1\pmod{3^{a+2}}\), so the first term in
\eqref{eq:shift-identity} vanishes modulo \(3^{a+1}\).  Substituting the two congruences for
\(S_L\) and \(T_L\) into \eqref{eq:shift-identity} proves \eqref{eq:3-adic-block-shift}.
\end{proof}

\begin{corollary}[\(3\)-adic lifting]\label{cor:3-adic-lifting}
Let \(R=2^\alpha5^\beta\), suppose \(R\mid k\) and \(3^a\mid P(k)\) with \(a\ge1\), and put
\begin{equation}\label{eq:3-adic-step}
L=\operatorname{lcm}(2\cdot3^a,R)=2c3^a,
\qquad 3\nmid c.
\end{equation}
Then exactly one of
\[
k,\qquad k+L,\qquad k+2L
\]
is a root modulo \(3^{a+1}\), while all three endpoints remain divisible by \(R\).  Moreover
\begin{equation}\label{eq:3-adic-period}
2\cdot3^e
\end{equation}
is a state period modulo \(3^e\).
\end{corollary}

\begin{proof}
The preceding lemma gives
\[
\frac{P(k+L)-P(k)}{3^a}\equiv c(2k+1)\pmod3.
\]
Since \(3\mid P(k)\) and \(10\equiv1\pmod3\),
\[
0\equiv P(k)\equiv1+2+\cdots+k=\frac{k(k+1)}2\pmod3,
\]
so \(k\equiv0\) or \(2\pmod3\).  Hence \(2k+1\not\equiv0\pmod3\), and the three
choices give the three possible residues at the next \(3\)-adic level.  Because \(R\mid L\),
the endpoint condition modulo \(R\) is preserved.

Taking \(L=2\cdot3^e\) in Lemma~\ref{lem:3-adic-shift} gives
\(P(k+L)\equiv P(k)\pmod{3^e}\) and \(k+L\equiv k\pmod{3^e}\), proving the period assertion.
\end{proof}

\section{A common prefix for every modulus}
\label{sec:common-prefix}

\begin{theorem}[A075002]\label{thm:A075002}
For every positive integer \(n\), there exists \(k\ge1\) such that
\[
n\mid P(k).
\]
Consequently OEIS A075002 is well-defined for every \(n\).
\end{theorem}

\begin{proof}
Factor
\begin{equation}\label{eq:A075002-factorization}
n=R\,3^e\prod_{i=1}^r p_i^{e_i},
\qquad
R=2^\alpha5^\beta,
\qquad
7\le p_1<\cdots<p_r.
\end{equation}
Let
\[
D=\operatorname{lcm}_{1\le i\le r}\ord_{p_i^{e_i}}(10),
\]
with \(D=1\) if \(r=0\).  We shall work in one sufficiently long \(d\)-digit block, where
\begin{equation}\label{eq:A075002-digit-choice}
d\equiv1\pmod D,
\qquad d\ge\max(\alpha,\beta).
\end{equation}
Such \(d\) are arbitrarily large.  Put \(B=10^d\).  Then for every \(i\),
\[
B\equiv10\pmod{p_i^{e_i}}.
\]

We first make the capacity choice precise.  For fixed \(n\), every step length appearing below
is obtained by taking least common multiples of \(R\), prime powers dividing \(n\), and
multiplicative orders modulo those prime powers.  Hence all possible step lengths, and the
maximum total displacement obtained by making at most \(p-1\) choices at each \(p\)-adic stage,
are bounded by a finite constant \(W(n)\) independent of \(d\).  Let
\begin{equation}\label{eq:A075002-initial-modulus}
M_0=R\cdot3^{\min(e,1)}.
\end{equation}
Choose \(d\) in the congruence class in \eqref{eq:A075002-digit-choice} so large that
\begin{equation}\label{eq:A075002-capacity}
9\cdot10^{d-1}>M_0+W(n).
\end{equation}
Let \(k_0\) be the least multiple of \(M_0\) that is at least \(10^{d-1}\).  Then
\(0\le k_0-10^{d-1}<M_0\), and \eqref{eq:A075002-capacity} keeps the entire construction
below inside this single \(d\)-digit block.

Since \(R\mid k_0\) and \(B\equiv0\pmod R\),
\begin{equation}\label{eq:A075002-initial-R}
P(k_0)=BP(k_0-1)+k_0\equiv0\pmod R.
\end{equation}
If \(e>0\), then \(3\mid k_0\), and because \(10\equiv1\pmod3\),
\[
P(k_0)\equiv\frac{k_0(k_0+1)}2\equiv0\pmod3.
\]

If \(e>1\), repeatedly apply Corollary~\ref{cor:3-adic-lifting}.  At stage \(a\), use
\[
L_a=\operatorname{lcm}(2\cdot3^a,R).
\]
One of the three choices raises a root modulo \(3^a\) to a root modulo \(3^{a+1}\), and the
\(R\)-condition is preserved.  Thus we obtain an endpoint \(k\) with
\[
R3^e\mid P(k).
\]
If \(e>0\), set \(H=2\cdot3^e\); otherwise set \(H=1\).  In either case \(H\) is a common
state period for all already installed odd prime-power conditions.

Now process \(p_1,\ldots,p_r\) in increasing order.  Suppose all preceding prime powers have
been installed at the current endpoint \(k\), and \(H\) is a common state period for them.  Let
\(p=p_i\) and
\[
h=\ord_p(10).
\]
Since \(B\equiv10\pmod p\), we also have \(h=\ord_p(B)\), and
\(p\nmid B(B-1)\).

Set
\begin{equation}\label{eq:A075002-first-step}
J_0=\operatorname{lcm}(h,H,R).
\end{equation}
Every prime factor of the period contributed by an earlier prime \(q<p\) is at most \(q\),
while \(R\) has no prime divisors other than \(2\) and \(5\).  Therefore \(p\nmid HR\), and also \(p\nmid h\) since
\(h\mid p-1\).  Hence \(J_0=c_0h\) with \(p\nmid c_0\).  Lemma~\ref{lem:unit-translation}
with \(a=0\) shows that
\[
P(k+tJ_0),\qquad 0\le t<p,
\]
runs through all residues modulo \(p\).  Because \(J_0\) is a multiple of \(H\) and \(R\),
every previously installed congruence and the endpoint condition modulo \(R\) are preserved.
Choose the unique \(t\) giving a root modulo \(p\).

Suppose inductively that a root modulo \(p^a\) has been obtained, where \(1\le a<e_i\).  Set
\begin{equation}\label{eq:A075002-lift-step}
J_a=\operatorname{lcm}(hp^a,H,R).
\end{equation}
Again \(p\nmid HRh\), so \(J_a=c_ahp^a\) with \(p\nmid c_a\).  Lemma~\ref{lem:unit-translation}
gives
\[
P(k+J_a)-P(k)\equiv c_ap^aU\pmod{p^{a+1}},
\qquad p\nmid U.
\]
Thus exactly one of
\[
k,\ k+J_a,\ldots,k+(p-1)J_a
\]
is a root modulo \(p^{a+1}\), while all earlier congruences remain intact.  Repeating installs
\(p^{e_i}\).

After completing this prime power, replace
\begin{equation}\label{eq:A075002-period-update}
H\leftarrow
\operatorname{lcm}\!\left(
H,
p^{e_i}\ord_{p^{e_i}}(B)
\right).
\end{equation}
By Lemma~\ref{lem:state-period}, the new factor is a state period modulo \(p^{e_i}\).  Since
\(B\equiv10\pmod{p_i^{e_i}}\), it equals
\(p^{e_i}\ord_{p^{e_i}}(10)\).  Every prime divisor of this factor is at most \(p\), so the
same argument applies at the next, larger prime.

At termination, the same endpoint \(k\) satisfies every prime-power condition in the
factorization \eqref{eq:A075002-factorization}.  Therefore \(n\mid P(k)\).
\end{proof}

\begin{remark}\label{rem:increasing-primes}
The increasing-prime order is essential to this proof.  A single global affine orbit need not
have unit translation simultaneously at all prime factors.  The construction above avoids that
obstruction by preserving completed states with periods whose prime divisors are all at most the
prime just processed; a larger, unprocessed prime therefore cannot divide the next step's
coefficient.
\end{remark}

\section{Verification and provenance}
\label{sec:verification}

The proofs are independent of computation.  A separate verifier evaluates prefixes modulo the
requested modulus using fixed-block transition matrices and cross-checks each matrix jump
against a literal recurrence on small instances.  The recorded checks include the published
initial values of A075002, the A075001 construction for \(n\le2000\), both lifting lemmas, and
common-index constructions for \(n\le500\), together with the five composite moduli below
\(5000\) for which a single affine orbit cannot provide unit translations at all prime factors.  These computations are
regression and reproducibility checks, not premises of the proofs.

The source, bibliography, verifier, verification record, and independent audit are preserved in
the public repository \href{https://github.com/dongxuelian2/oeis-papers}{\texttt{oeis-papers}}.
The frozen release \href{https://github.com/dongxuelian2/oeis-papers/releases/tag/a075000-a075002-v1.0.1}{\texttt{a075000-a075002-v1.0.1}}
contains this paper as a stable snapshot.  From the paper directory, the verifier can be
reproduced with
\begin{verbatim}
python verify_concatenation_divisibility.py
\end{verbatim}
The manuscript grew out of a larger computer-assisted investigation, but that research history
is not used as a premise here.  The OEIS status statement above records only the checked date;
it is not intended as a claim of literature-wide priority.

\bibliographystyle{amsplain}
\bibliography{references}

\end{document}
